\documentclass[main.tex]{subfiles}


\begin{document}

%from pagenumber to up
% \phantomsection 
% \hypertarget{toc}{}

 

 

% номер секции вручную
\setcounter{section}{21
}
\addtocounter{section}{-1} 

\section{Работа тока. Закон Джоуля"--~Ленца}


Электрический ток несет с собой \emph{энергию} (аналогия: поток жидкости, несущий механичесую энергию). 
% Энергия, переносимая током, может быть превращена в любую другую форму (тепло, излучение и~т.\,д.). 
Эта энергия (\emph{электроэнергия}) возникает за счет работы электрического поля, которое <<гонит>> заряды тока по проводнику. 

\textbf{Работа тока} ($A$ [Дж])~--- это работа электрического поля по передвижению зарядов тока в проводнике:
\begin{equation}
    A=UIt,
\end{equation}
где $U$ и~$I$~--- напряжение и ток проводника, $t$~--- время протекания тока.

\textbf{Мощность тока} ($P$ [Вт])~--- это быстрота совершения работы тока: 
\begin{equation}
    P=UI.
\end{equation}
% Мощность, потребляемая электроприбором,~--- это и есть мощность тока в нем.

\textbf{Электроэнергия}, потребленная прибором, может быть рассчитана так:
\begin{equation}
    W=Pt.
\end{equation}

Практически всегда возникает необходимость преобразовать электроэнергию в другой вид энергии. В связи с этим широкое применение нашли следующие устройства.
\begin{itemize}
    \item \emph{Электронагреватель} преобразует электроэнергию в теплоту.

    \textbf{КПД электронагревателя} равен:
    \begin{equation}
        \eta=\frac{Q_\text{п}}{W_\text{з}},
    \end{equation}
    где $Q_\text{п}$~--- полезное тепло, $W_\text{з}$~--- затраченная электроэнергия.

    \item \emph{Электродвигатель} преобразует электроэнергию в работу.

    \textbf{КПД электродвигателя} равен:
    \begin{equation}
        \eta=\frac{A_\text{п}}{W_\text{з}},
    \end{equation}
    где $A_\text{п}$~--- полезная работа, $W_\text{з}$~--- затраченная электроэнергия.    
\end{itemize}

Пусть к батарейке подключен стальной провод (рис.\,\ref{fig:p125}).


\tikzfading[name=fade out,
inner color=transparent!0,
outer color=transparent!100]


    \begin{figure}[H]
    \centering
\begin{circuitikz}[font=\small,thin,
    line cap=round,line join=round,
>={Stealth[inset=0pt,round,length=3mm, angle'=20]},
vec/.style={>={Stealth[inset=0pt,round,length=3.5mm, angle'=20]}}]

    \ctikzset{bipoles/americaninductor/coils=3}


\path[] (1.9,.75) --++ (.2,0) --++ (0,-.25) coordinate (r) --++ (0,-.25) --++ (-.2,0) -- cycle; 


%cond
\path (r) --++ (.5,0) --++ (0,-.75) coordinate (I) --++ (0,-.75) --++ (-.6,0) coordinate (rr) ++ (-1,0) coordinate (R) ++ (-1,0) --++ (-.6,0) --++ (0,1.5) --++ (.5,0);
\fill[red,semitransparent] ([yshift=-.25cm]rr) rectangle++ (-2,.5);
\draw[] ([yshift=-.25cm]rr) rectangle++ (-2,.5);
\draw (r) --++ (.5,0) --++ (0,-1.5) --++ (-.6,0) coordinate (rr) ++ (-2,0) --++ (-.6,0) --++ (0,1.5) --++ (.5,0);
\node[below] at ([yshift=-.25cm]R) {$R$};


%I
\draw[->,red] ([shift={(.2,.5)}]I) --++ (0,-1) node[midway, right, black] {$I$}; 


%terminals
\begin{scope}
    \clip[] (2,0) rectangle++ (.2,1);
    \filldraw[gray,rounded corners=1pt] (1.9,.75) --++ (.2,0) --++ (0,-.25) coordinate (r) --++ (0,-.25) --++ (-.2,0) -- cycle; 
\end{scope}

\begin{scope}
    \clip[] (0,0) rectangle++ (-.2,1);
    \filldraw[gray,rounded corners=1pt] (.1,.85) --++ (-.2,0) --++ (0,-.35) coordinate (l) --++ (0,-.35) --++ (.2,0) -- cycle; 
\end{scope}


%batt
\begin{scope}
    \clip[rounded corners=1mm] (0,0) rectangle++ (2,1);
    \fill[NavyBlue,nearly transparent] (0,0) rectangle++ (1,1);
    \fill[red,nearly transparent] (1,0) rectangle++ (1,1);
\end{scope}
\draw[rounded corners=1mm] (0,0) rectangle++ (2,1);


%labels
\node[right] at (0,.5) {$-$};
\node[left] at (2,.5) {$+$};
\node[above] at (1,1) {Б}; 


% %heat
\path ([yshift=-.25cm]rr) coordinate (rec);
% \foreach \i in {1,...,100}
% \draw[red,opacity={1-\i/100},line width=.1pt] ([shift={(\i*.9*\pgflinewidth,-\i*.9*\pgflinewidth)}]rec) rectangle++ ({-2cm-2*\i*.9*\pgflinewidth},{.5cm+2*\i*.9*\pgflinewidth}); 


% \fill[path fading=south,red] (rec) rectangle++ (-2,-.1);
% \begin{scope}
%     \clip ([shift={(.0001,0)}]rec) rectangle++ (.1,-.1);
%     \fill[path fading=fade out,red] ([shift={(-.1,.1)}]rec) rectangle++ (.2,-.2);
% \end{scope}



\end{circuitikz}
\caption{Опыт с батарейкой и проводом}
\label{fig:p125}
\end{figure}

Батарейка~Б вызывает протекание тока\footnote{Направлением тока принято считать направление движения \emph{положительных} зарядов (в данном случае они как бы вытекают из <<плюса>> батарейки и втекают в ее <<минус>>.)
}~$I$ в проводе с сопротивлением~$R$. Опыт показывает, что за время~$t$ провод \emph{нагреется}~--- в нем выделится \emph{тепло}, которое можно вычислить по \textbf{закону Джоуля"--~Ленца}:
\begin{equation}
    Q=I^2Rt.
\end{equation}


% \emph{Тепло}, выделяющееся на любом участке цепи при протекании по нему тока, можно вычислить по \textbf{закону Джоуля"--~Ленца}:
% \begin{equation}
    % Q=I^2Rt
% \end{equation}
% где $I$ и~$R$~--- ток и сопротивление участка цепи, $t$~--- время протекания тока.





 
\end{document}